% Generated draft: requires mathematical and source review.
\documentclass[11pt]{article}
\usepackage[margin=25mm]{geometry}
\usepackage{fontspec}
\setmainfont{Latin Modern Roman}
\usepackage{amsmath,amssymb,mathtools,unicode-math}
\setmathfont{Latin Modern Math}
\usepackage{enumitem,needspace,xurl,hyperref}
\setlength{\parindent}{0pt}
\setlength{\parskip}{5pt}
\setlength{\emergencystretch}{3em}
\newcommand{\N}{\mathbb{N}}
\newcommand{\Z}{\mathbb{Z}}
\newcommand{\Q}{\mathbb{Q}}
\newcommand{\R}{\mathbb{R}}
\newcommand{\C}{\mathbb{C}}
\newcommand{\F}{\mathbb{F}}
\newcommand{\A}{\mathbb{A}}
\newcommand{\PP}{\mathbb{P}}
\newcommand{\E}{\mathbb{E}}
\newcommand{\eps}{\varepsilon}
\newcommand{\dd}{\,\mathrm{d}}
\begin{document}
\section*{UnsolvedMath ID: 30004761}
\textbf{Source problem code:} catalogue entry
\subsection{Definitions and mathematical statement}
For positive integers \(m,n\), a function \(f\colon\R^{m\times n}\to\R\) is rank-one convex if
\[
f((1-t)A+tB)\leq (1-t)f(A)+tf(B)
\]
for all \(A,B\in\R^{m\times n}\) with \(\operatorname{rank}(A-B)\leq1\) and all \(t\in[0,1]\). This gives the standard full-matrix-space definition; the supplied record does not specify the dimensions, the domain of the integrand, or its regularity.

To explain outer variations conditionally, suppose that \(\Omega\subset\R^n\) is open, \(f\) is continuously differentiable, and \(u\colon\Omega\to\R^m\) has sufficient regularity and integrability for the following expressions and differentiation to be valid. For the energy
\[
I_f(u)=\int_\Omega f(Du(x))\,\dd x,
\]
outer stationarity means
\[
\left.\frac{\dd}{\dd t}I_f(u+t\varphi)\right|_{t=0}
=\int_\Omega Df(Du):D\varphi\,\dd x=0
\qquad\text{for every }\varphi\in C_c^\infty(\Omega;\R^m).
\]
Here \(A:B=\sum_{i,j}A_{ij}B_{ij}\); equivalently, \(\operatorname{div}Df(Du)=0\) in distributions. These are variations of the dependent variable. No stationarity under changes of the independent variable is imposed. This is a conditional explanation of the terminology, not a specification of the missing hypotheses in the source problem.

A parametrized Young measure on a finite-dimensional phase space \(E\) is a weakly measurable family \((\nu_x)_{x\in\Omega}\) of Borel probability measures on \(E\). Call it trivial if \(\nu_x\) is a Dirac probability measure for almost every \(x\in\Omega\). For a homogeneous Young measure, triviality means that its single probability measure is a Dirac mass. The supplied record does not identify the phase space, whether the measures are homogeneous, or which generating sequences and differential constraints are admissible.

\textbf{Requested mathematical target, with unresolved qualifications.} Find rank-one convex integrands \(f\) for which every admissible Young measure associated with the differential inclusion for critical points under outer variations only is trivial. If \(\mathcal Y_{\mathrm{out}}(f)\) denotes the intended admissible class, the desired property is
\[
\forall\nu\in\mathcal Y_{\mathrm{out}}(f),\qquad
\nu_x\text{ is a Dirac mass for almost every }x.
\]
The notation \(\mathcal Y_{\mathrm{out}}(f)\) is a placeholder for information absent from the supplied record. In particular, the actual differential inclusion, the admissibility and generation conditions, growth assumptions, boundary conditions, and dimensions remain unspecified. Consequently, the record does not determine a unique fully quantified problem. The quantifier over Young measures must not be interpreted as ranging over all probability measures supported on an unspecified inclusion set.
\subsection{Short English statement}
Find rank-one convex energies for which the differential inclusion describing critical points under outer variations admits only Dirac Young measures. The supplied record leaves the inclusion and admissible class of measures undefined, so a precise final formulation requires mathematical source review. Its partially solved status and claimed partial results are supplied dated curation, not independently verified conclusions.
\subsection{Sources}
\begin{itemize}
\item Supplied local problem registry record, permanent UnsolvedMath ID \(30004761\), source problem code \texttt{catalogue entry}, title \emph{Rank-One Convex Energies Excluding Nontrivial Young Measures}, in the Oberwolfach Reports Open Problems collection. The record gives proposed year \(2021\) and identifies no proposer.
\item Supplied report citation: \emph{Partial Differential Equations (2022)}, Oberwolfach Reports, DOI \texttt{10.4171/owr/2021/35}, pp.\ 1859--1914. \url{https://doi.org/10.4171/owr/2021/35}. The supplied background describes Oberwolfach Reports as documenting mathematical workshops and their open-problem discussions. The citation's year \(2022\), the proposed year \(2021\), and the DOI's \texttt{2021} component are retained as supplied.
\item Supplied literature source labelled De Rosa--Tione: \url{https://arxiv.org/abs/2011.09922}. The record supplies no article title or precise theorem identifying its application to this problem.
\item Supplied literature source labelled Oberwolfach report: \url{https://ems.press/content/serial-article-files/46911}.
\item Literature assessment checked at \texttt{2026-08-22}: supplied dated curation, not independently verified. The supplied status is \texttt{partially\_solved}. The supplied assessment states that atomic-condition results exclude nontrivial Young measures for broad classes of stationary graphs, providing a partial answer, while the outer-variation-only question in full generality remains unresolved. The background expressly calls this dated literature triage derived from the supplied OWR curation and requests independent verification of time-sensitive status and cited claims.
\item The registry labels the statement \texttt{corrected\_verified} and reports that the official Oberwolfach report was checked directly, with the clean statement preserving the mathematical target while restoring definitions and removing extraction artifacts or source-dependent wording. This is a supplied verification claim; no independent source inspection is asserted here.
\end{itemize}
\end{document}
